\section{The Page You Cannot Ask For}
\label{sec:ch13-the-page-you-cannot-ask-for}
\index{pagination!sorting by a field somebody else owns|(}%

Show a client the graph Mosaic now serves and one of the first things they will
try is cheapest first.

\begin{minted}[fontsize=\footnotesize,breaklines]{graphql}
{ browseProducts(first: 3, order: [{ price: { amount: ASC } }]) { nodes { title } } }
\end{minted}

\begin{minted}[fontsize=\footnotesize,breaklines]{text}
Field "price" is not defined by type "ProductSortInput".
\end{minted}

Filtering goes the same way:

\begin{minted}[fontsize=\footnotesize,breaklines]{text}
Field "availableQuantity" is not defined by type "ProductFilterInput".
\end{minted}

Both are validation failures at the router, before any subgraph is asked
anything. And the composed inputs are exactly Catalog's own columns:

\begin{minted}[fontsize=\footnotesize,breaklines]{text}
sort:   id, sku, title, description, category
filter: and, or, id, sku, title, description, category
\end{minted}

\index{ProductSortInput@\code{ProductSortInput}!is Catalog's table}%
Which is correct, and is the shape of the problem rather than a gap in it. The
sort input is generated from the entity Catalog holds. Pricing contributes a
field to \code{Product}, and contributing a field to a type is not the same as
contributing a column to somebody else's table. The order of a page is decided
by whoever runs the query that produces the page, and only Catalog runs that
query.

\subsection*{Why no directive rescues this}

\index{requires@\code{@requires}!nested field sets}%
The tempting move is \code{@requires}. Catalog needs a price to sort by; there is
a directive whose entire job is getting a value from another subgraph into a
resolver; connect the two. Chapter~\ref{ch:entity-resolution-done-right} left two
questions about how far that directive stretches, so before arguing about
whether it fits, here is how far it does stretch.

A nested field set composes. Inventory can declare Pricing's \code{price} as
\code{@external} and require a path through it:

\begin{graphqlsdl}[fontsize=\footnotesize]
type Product @key(fields: "id") {
  availableQuantity: Int!
  restockThreshold: Int! @requires(fields: "price { amount }")
  id: ID!
  price: Money! @external
}
\end{graphqlsdl}

\index{requires@\code{@requires}!spanning two other subgraphs}%
And a field set spanning two other subgraphs composes. Pricing's
\code{shippingCost} can require \code{category}, which Catalog owns, and
\code{availableQuantity}, which Inventory owns, in one directive, and the routing
table afterwards has the three fields in three different services:

\begin{minted}[fontsize=\footnotesize]{text}
Product.category           catalog
Product.shippingCost       pricing
Product.availableQuantity  inventory
\end{minted}

So the answer to both of chapter~\ref{ch:entity-resolution-done-right}'s
questions is yes, and neither helps here. \code{@requires} feeds a field
resolver on an entity, and it runs after the router has located that entity. The
router locates an entity by being handed a key, and a key comes from the fetch
that produced the page. There is no page yet when the ordering has to be decided,
so there is nothing for the value to be required onto. The directive is not too
weak. It is in the wrong half of the request.

The same argument disposes of \code{@provides}, of an \code{@external} copy of
the price, and of every other federation directive, because all of them describe
what happens once the router knows which entities it is dealing with.

\subsection*{Three things that do work, and what each costs}

\index{pagination!denormalising a sort key}%
The first is to copy the sort key into the service that owns the page. Catalog
grows a \code{price\_cents} column, keeps it up to date from Pricing's events,
and sorts on it. This is the one most teams end up with and it is worth being
honest about what it is: a second copy of somebody else's data, with a staleness
window, in the service whose whole reason for existing was not knowing about
prices. What it buys is a page that sorts correctly in one query, and a page is
the thing clients ask for constantly.

\index{pagination!a service that owns the ordering}%
The second is to move the ordering somewhere that can see everything. A search
or listing service subscribes to both domains, holds an index, and answers
\code{searchProducts} with a page of \code{Product} keys the router then resolves
in the ordinary way. The page is one service's answer again, which is the
property section~\ref{sec:ch13-the-page-that-fits} showed is what makes a
connection work at all. The cost is a service, an index, and a second definition
of what a product is; the benefit is that this is also where full-text search,
facets and relevance were going to end up anyway.

The third is to admit the page is not sortable that way and give the client a
narrower field: \code{cheapestInCategory}, answered by Pricing, returning a
handful of keys. Small, honest, and it stops being enough the moment somebody
wants page four.

\index{federation!what it does not remove}%
What none of them is is a schema change. That is the thing to take away, and it
is worth stating plainly because a graph that presents six services as one
schema invites the assumption that it presents six databases as one database.
It does not. A query that ranks rows by a column two services would have to
share is a distributed join, and this stack has no distributed join in it. What
it has is a mechanism for fetching fields of things it has already found, which
is a different and much cheaper thing.

\medskip
That is a limit with no way round it. The next one turned out to have a way
round it, which cost a service.

\index{pagination!sorting by a field somebody else owns|)}%
